% Pravila sklepanja iz poglavja Dokazi na enem listu A4.
\documentclass[10pt]{article}

\usepackage[a4paper, landscape, margin=10mm]{geometry}
\usepackage[slovene]{babel}

\usepackage{amsmath}
\usepackage{mathtools}

\usepackage{fontspec}
\usepackage{unicode-math}
\setmainfont{Libertinus Serif}[Ligatures = TeX, Numbers = OldStyle]
\setmathfont{Libertinus Math}
\newcommand{\lmnVerzalke}{\addfontfeature{Numbers=Lining}}
\usepackage[protrusion = true, expansion = true, final]{microtype}

\usepackage{xcolor}
\definecolor{ULantracit}{HTML}{58595B}
\definecolor{ULsiva}{HTML}{A7A8AA}
\colorlet{barvaPasu}{ULsiva!35}

\usepackage{multicol}
\setlength{\columnsep}{6mm}

\newcommand{\nls}[1]{\guillemotright\textit{#1}\guillemotleft} % naravni jezik
\newcommand{\df}[1]{\emph{\textbf{#1}}}  % definiranec
\newcommand{\defeq}{\mathrel{{:}{=}}}  % definicijska enacba
\newcommand{\dfeq}{\defeq}                     % starejse ime

\newcommand{\konst}[1]{\mathrm{k}_{#1}} % konstantna preslikava


% Substitucija
\newcommand{\xsubst}[2]{#1[#2]}
\newcommand{\subst}[2]{(#1)[#2]}
\newcommand{\subto}{\,{\mapsto}\,}

%%% quantifiers
\newcommand{\all}[1]{\forall #1 \,.\,}
\newcommand{\some}[1]{\exists #1 \,.\,}
\newcommand{\exactlyone}[1]{\exists! #1 \,.\,}
\newcommand{\descr}[1]{\iota #1 \,.\,}

\newcommand{\defiff}{\quad {{:}{\liff}}\quad }

% \Box je v unicode-math prazen kvadratek; rabimo ga kot vsebnik.
\renewcommand{\Box}{\mdlgwhtsquare}

\newcommand{\brac}[1]{\ensuremath{\cdots}\enspace$\langle$#1$\rangle$\enspace\ensuremath{\cdots}}

%\newcommand{\lam}[3]{\lambda #1 \,{\in}\, #2\,.\left(#3\right)}

% \newcommand{\uall}[2]{\forall\, #1\,.\left(#2\right)}
% \newcommand{\usome}[2]{\exists\, #1\,.\left(#2\right)}
% \newcommand{\uexactlyone}[3]{\exists!\, #1\,.\left(#2\right)}
% \newcommand{\ulam}[2]{\lambda #1 .\left(#2\right)}

% \newcommand{\xall}[3]{\forall\, #1 \,{\in}\, #2\,.\,#3}
% \newcommand{\xsome}[3]{\exists\, #1 \,{\in}\, #2\,.\,#3}
% \newcommand{\xexactlyone}[3]{\exists!\, #1 \,{\in}\, #2\,.\,#3}
% \newcommand{\xlam}[3]{\lambda #1 \,{\in}\, #2\,.\,#3}

% \newcommand{\xuall}[2]{\forall\, #1\,.\,#2}
% \newcommand{\xusome}[2]{\exists\, #1\,.\,#2}
% \newcommand{\xuexactlyone}[2]{\exists!\, #1,.\,#2}
% \newcommand{\xulam}[2]{\lambda #1 .\,#2}

% \newcommand{\tlam}[3]{\lambda #1 \,{:}\, #2\,.\,\left(#3\right)}
% \newcommand{\xtlam}[3]{\lambda #1 \,{:}\, #2\,.\,#3}

\newcommand{\iinn}{\qquad\text{in}\qquad}

% Implikacija
\newcommand{\lthen}{\Rightarrow}
\newcommand{\liff}{\Leftrightarrow}

% Hipoteticna izpeljava: iz predpostavk #1 pridemo do #2.
\NewDocumentCommand{\izpeljava}{m m}{\begin{array}[b]{c}{#1}\\ \vdots\\ {#2}\end{array}}
\newcommand{\hypoinfer}{\izpeljava}  % starejse ime

% Standardne množice
\newcommand{\NN}{\mathbb{N}}
\newcommand{\ZZ}{\mathbb{Z}}
\newcommand{\QQ}{\mathbb{Q}}
\newcommand{\RR}{\mathbb{R}}
\newcommand{\CC}{\mathbb{C}}

\newcommand{\one}{\mathbb{1}}
\newcommand{\two}{\mathbb{2}}
\newcommand{\unit}{(\,)}

%%% Množice
\newcommand{\set}[1]{\left\{#1\right\}}
\newcommand{\such}{\; \middle| \;}
\newcommand{\iso}{\cong}

%%%%%%  Projections and Injections
%%%%%%%%%%%%%%%%%%%%%%%%%%%%%%%%%%%%%%%%%%%%%%%%%%%%%%%%%%%%%
\newcommand{\inj}[1][]{\mathsf{in}_{#1}}
\newcommand{\pr}[1][]{\mathsf{pr}_{#1}}

%% Neobvezna nadpisa loci vejica, kadar sta podana oba.
\NewDocumentCommand{\nadpisa}{m m}{#1\IfBlankF{#2}{,}#2}

\NewDocumentCommand{\fst}{O{} O{}}{\pr[1]^{\nadpisa{#1}{#2}}}
\NewDocumentCommand{\snd}{O{} O{}}{\pr[2]^{\nadpisa{#1}{#2}}}
\NewDocumentCommand{\inl}{O{} O{}}{\inj[1]^{\nadpisa{#1}{#2}}}
\NewDocumentCommand{\inr}{O{} O{}}{\inj[2]^{\nadpisa{#1}{#2}}}


%%% Naravna števila
\newcommand{\suc}[1]{{#1^{+}}}

\newcommand{\Set}{\mathsf{Set}}
\newcommand{\On}{\mathsf{On}}

%% Drevesa
\newcommand{\Tree}{\mathsf{Tree}}
\newcommand{\tree}[1]{\mathsf{tree}(#1)}
\newcommand{\emptyTree}{\mathsf{empty}}

%%% Operacije na množicah

\newcommand{\pow}[1]{\mathcal{P}(#1)}       % potenčna množica
\newcommand{\npow}[2]{\mathcal{P}_{#1}(#2)} % podmnožice dane velikosti

%% Preslikave
\newcommand{\id}[1][]{\mathrm{id}_{#1}}
\newcommand{\inv}[1]{{#1}^{-1}}
\newcommand{\restrict}[2]{#1\mathrel{\big|}_{#2}}

\newcommand{\diag}[1][]{\Delta_{#1}}

\newcommand{\img}[1]{#1_{*}}
\newcommand{\invimg}[1]{#1^{*}}

%% Kombinatorika
\newcommand{\ceil}[1]{\lceil #1 \rceil}

%% Relacije
\newcommand{\transpose}[1]{#1^{\mathsf{T}}}
\newcommand{\compl}[1]{#1^{\mathsf{C}}}
\newcommand{\rel}[1]{\mathrel{#1}}


\usepackage{boxedproof}
\renewcommand{\boxNumberStyle}{\scriptsize\lmnVerzalke}
\colorlet{boxProofColor}{ULantracit}
\renewcommand{\boxProofSize}{\footnotesize}

\pagestyle{empty}
\setlength{\parindent}{0pt}

%% Poravnava: vsak zunanji okvir je natanko \sirinaOkvirja sirok, vgnezdeni
%% okvirji pa segajo do enakega odmika od levega in desnega roba.
\newlength{\sirinaOkvirja}
\newlength{\razmikPara}  \setlength{\razmikPara}{4mm}
\newcommand{\polna}{\setlength{\sirinaOkvirja}{\linewidth}}
\newcommand{\polovica}{%
  \setlength{\sirinaOkvirja}{\dimexpr(\linewidth - \razmikPara)/2\relax}}

\makeatletter
\renewenvironment{boxFrame}{%
  \setbox\box@content=\hbox\bgroup
    \begin{varwidth}{\linewidth}\setlength{\parindent}{0pt}}
 {\end{varwidth}\egroup
  \setbox\box@content=\hbox to
    \dimexpr\sirinaOkvirja - 2\fboxsep - 2\fboxrule\relax
    {\box\box@content\hfil}%
  \fcolorbox{boxProofColor}{white}{\box\box@content}%
  \ignorespacesafterend}
\renewcommand{\box@computeTarget}{%
  \setlength{\boxTargetWidth}{%
    \dimexpr\sirinaOkvirja - 4\fboxsep - 4\fboxrule - 2\boxThisIndent\relax}}
\makeatother

%% Razdelek: pas z imenom in znakom, nato pari oznak in vzorcev.
\newlength{\razmikRazdelkov} \setlength{\razmikRazdelkov}{2ex}

\newcommand{\pas}[2]{%
  {\setlength{\fboxsep}{1.5pt}%
   \colorbox{barvaPasu}{\makebox[\dimexpr\linewidth - 2\fboxsep\relax][l]{%
       \strut\hspace{1pt}\textbf{#1}\hfill\ensuremath{#2}\hspace{1pt}}}}}

\newenvironment{razdelek}[2]{%
  \par\noindent\begin{minipage}{\linewidth}%
    \polna
    \pas{#1}{#2}\par\vspace{1ex}}
 {\end{minipage}\par\vspace{\razmikRazdelkov}}

\newcommand{\oznaka}[1]{%
  \par\addvspace{1.2ex}\polna
  {\footnotesize\color{ULantracit}\textsc{#1}}\par\nopagebreak\vspace{0.3ex}}

\newcommand{\opomba}[1]{\par\vspace{0.3ex}{\footnotesize\itshape #1}\par}

\begin{document}

\raggedcolumns
\begin{multicols*}{3}

{\Large\bfseries Pravila sklepanja}\par
\vspace{0.3ex}
{\footnotesize\color{ULantracit}\textsc{Logika in množice}}\par
\vspace{\razmikRazdelkov}

\begin{razdelek}{Konjunkcija}{\land}
  \oznaka{vpeljava}
  \begin{boxedProof}
    \begin{boxFrame}
      \boxLine{Dokažimo $\phi \land \psi$:}
      \begin{boxNested}
        \boxLine{Dokažimo $\phi$:}
        \boxLine[1]{\brac{dokaz $\phi$}}
      \end{boxNested}
      \begin{boxNested}
        \boxLine{Dokažimo $\psi$:}
        \boxLine[1]{\brac{dokaz $\psi$}}
      \end{boxNested}
    \end{boxFrame}
  \end{boxedProof}
  \oznaka{uporaba}
  \polovica
  \begin{boxedProof}[reserved]%
    \begin{boxFrame}
      \boxLineTag{$k$}{Velja $\phi \land \psi$.}
      \boxGap
      \boxLine{Zaradi~($k$) velja $\phi$.}
    \end{boxFrame}
  \end{boxedProof}%
  \hspace{\razmikPara}%
  \begin{boxedProof}[reserved]%
    \begin{boxFrame}
      \boxLineTag{$k$}{Velja $\phi \land \psi$.}
      \boxGap
      \boxLine{Zaradi~($k$) velja $\psi$.}
    \end{boxFrame}
  \end{boxedProof}
\end{razdelek}

\begin{razdelek}{Implikacija}{\lthen}
  \oznaka{vpeljava}
  \begin{boxedProof}
    \begin{boxFrame}
      \boxLine{Dokažimo $\phi \lthen \psi$:}
      \boxLine[1]{Predpostavimo $\phi$.}
      \boxLine[1]{Dokažimo $\psi$:}
      \boxLine[2]{\brac{dokaz $\psi$}}
    \end{boxFrame}
  \end{boxedProof}
  \oznaka{uporaba}
  \begin{boxedProof}[reserved]%
    \begin{boxFrame}
      \boxLineTag{$k$}{Velja $\phi \lthen \psi$.}
      \boxGap
      \boxLineTag{$m$}{Velja $\phi$.}
      \boxGap
      \boxLine{Zaradi~($k$) in~($m$) velja $\psi$.}
    \end{boxFrame}
  \end{boxedProof}
\end{razdelek}

\begin{razdelek}{Negacija}{\lnot}
  \oznaka{vpeljava}
  \begin{boxedProof}
    \begin{boxFrame}
      \boxLine{Dokažimo $\lnot \phi$:}
      \boxLine[1]{Predpostavimo $\phi$.}
      \boxLine[1]{Dokažimo~$\bot$:}
      \boxLine[2]{\brac{dokaz $\bot$}}
    \end{boxFrame}
  \end{boxedProof}
  \oznaka{uporaba}
  \begin{boxedProof}[reserved]%
    \begin{boxFrame}
      \boxLineTag{$k$}{Velja $\phi$.}
      \boxGap
      \boxLineTag{$m$}{Velja $\neg \phi$.}
      \boxGap
      \boxLine{Zaradi~($k$) in~($m$) velja $\bot$.}
    \end{boxFrame}
  \end{boxedProof}
\end{razdelek}

\begin{razdelek}{Disjunkcija}{\lor}
  \oznaka{vpeljava}
  \polovica
  \begin{boxedProof}
    \begin{boxFrame}
      \boxLine{Dokažimo $\phi \lor \psi$:}
      \boxLine[1]{Dokažimo levi disjunkt $\phi$:}
      \boxLine[2]{\brac{dokaz $\phi$}}
    \end{boxFrame}
  \end{boxedProof}%
  \hspace{\razmikPara}%
  \begin{boxedProof}
    \begin{boxFrame}
      \boxLine{Dokažimo $\phi \lor \psi$:}
      \boxLine[1]{Dokažimo desni disjunkt $\psi$:}
      \boxLine[2]{\brac{dokaz $\psi$}}
    \end{boxFrame}
  \end{boxedProof}
  \oznaka{uporaba}
  \begin{boxedProof}[reserved]%
    \begin{boxFrame}
      \boxLineTag{$k$}{Velja $\phi \lor \psi$.}
      \boxGap
      \boxLine{Dokažimo $\rho$:}
      \boxLine[1]{Uporabimo ($k$) in obravnavamo dva primera:}
      \begin{boxNested}[1]
        \boxLine{Predpostavimo $\phi$.}
        \boxLine{Dokažimo $\rho$:}
        \boxLine[1]{\brac{dokaz $\rho$}}
      \end{boxNested}
      \begin{boxNested}[1]
        \boxLine{Predpostavimo $\psi$.}
        \boxLine{Dokažimo $\rho$:}
        \boxLine[1]{\brac{dokaz $\rho$}}
      \end{boxNested}
    \end{boxFrame}
  \end{boxedProof}
\end{razdelek}

\begin{razdelek}{Resnica in neresnica}{\top\;\bot}
  \oznaka{vpeljava resnice}
  \begin{boxedProof}
    \begin{boxFrame}
      \boxLine{Dokažimo $\top$:}
      \boxLine[1]{Očitno.}
    \end{boxFrame}
  \end{boxedProof}
  \oznaka{uporaba neresnice}
  \begin{boxedProof}[reserved]%
    \begin{boxFrame}
      \boxLineTag{$k$}{Velja $\bot$.}
      \boxGap
      \boxLine{Zaradi~($k$) velja $\rho$.}
    \end{boxFrame}
  \end{boxedProof}
  \opomba{Resnica nima pravila uporabe, neresnica nima pravila vpeljave.}
\end{razdelek}

\begin{razdelek}{Ekvivalenca}{\liff}
  \oznaka{vpeljava}
  \begin{boxedProof}
    \begin{boxFrame}
      \boxLine{Dokažimo $\phi \liff \psi$:}
      \begin{boxNested}
        \boxLine{Dokažimo $\phi \lthen \psi$:}
        \boxLine[1]{\brac{dokaz $\phi \lthen \psi$}}
      \end{boxNested}
      \begin{boxNested}
        \boxLine{Dokažimo $\psi \lthen \phi$:}
        \boxLine[1]{\brac{dokaz $\psi \lthen \phi$}}
      \end{boxNested}
    \end{boxFrame}
  \end{boxedProof}
  \oznaka{uporaba}
  \polovica
  \begin{boxedProof}[reserved]%
    \begin{boxFrame}
      \boxLineTag{$k$}{Velja $\phi \liff \psi$.}
      \boxGap
      \boxLine{Zaradi~($k$) velja $\phi \lthen \psi$.}
    \end{boxFrame}
  \end{boxedProof}%
  \hspace{\razmikPara}%
  \begin{boxedProof}[reserved]%
    \begin{boxFrame}
      \boxLineTag{$k$}{Velja $\phi \liff \psi$.}
      \boxGap
      \boxLine{Zaradi~($k$) velja $\psi \lthen \phi$.}
    \end{boxFrame}
  \end{boxedProof}
\end{razdelek}

\begin{razdelek}{Univerzalni kvantifikator}{\forall}
  \oznaka{vpeljava}
  \begin{boxedProof}
    \begin{boxFrame}
      \boxLine{Dokažimo $\all{x \in A} \phi(x)$:}
      \boxLine[1]{Naj bo $x \in A$.}
      \boxLine[1]{Dokažimo $\phi(x)$:}
      \boxLine[2]{\brac{dokaz $\phi(x)$}}
    \end{boxFrame}
  \end{boxedProof}
  \opomba{Spremenljivka $x$ mora biti sveža.}
  \oznaka{uporaba}
  \begin{boxedProof}[reserved]%
    \begin{boxFrame}
      \boxLineTag{$k$}{Velja $\all{x \in A} \phi(x)$.}
      \boxGap
      \boxLineTag{$m$}{Velja $\langle\text{izraz}\rangle \in A$.}
      \boxGap
      \boxLine{Zaradi~($k$) in~($m$) velja $\phi(\langle\text{izraz}\rangle)$.}
    \end{boxFrame}
  \end{boxedProof}
\end{razdelek}

\begin{razdelek}{Eksistenčni kvantifikator}{\exists}
  \oznaka{vpeljava}
  \begin{boxedProof}
    \begin{boxFrame}
      \boxLine{Dokažimo $\some{x \in A} \phi(x)$:}
      \boxLine[1]{Podamo $x \defeq \langle\text{izraz}\rangle$.}
      \begin{boxNested}[1]
        \boxLine{Dokažimo $\langle\text{izraz}\rangle \in A$:}
        \boxLine[1]{\brac{dokaz $\langle\text{izraz}\rangle \in A$}}
      \end{boxNested}
      \begin{boxNested}[1]
        \boxLine{Dokažimo $\phi(\langle\text{izraz}\rangle)$:}
        \boxLine[1]{\brac{dokaz $\phi(\langle\text{izraz}\rangle)$}}
      \end{boxNested}
    \end{boxFrame}
  \end{boxedProof}
  \oznaka{uporaba}
  \begin{boxedProof}[reserved]%
    \begin{boxFrame}
      \boxLineTag{$k$}{Velja $\some{x \in A} \phi(x)$.}
      \boxGap
      \boxLine{Zaradi~($k$) imamo $x \in A$ in $\phi(x)$.}
    \end{boxFrame}
  \end{boxedProof}
  \opomba{Spremenljivka $x$ mora biti sveža.}
\end{razdelek}

\begin{razdelek}{Drugo}{}
  \oznaka{dokaz s protislovjem}
  \begin{boxedProof}
    \begin{boxFrame}
      \boxLine{Dokažimo $\phi$ s protislovjem:}
      \boxLine[1]{Predpostavimo $\neg\phi$.}
      \boxLine[1]{Dokažimo~$\bot$:}
      \boxLine[2]{\brac{dokaz $\bot$}}
    \end{boxFrame}
  \end{boxedProof}
  \oznaka{izpeljava dejstva}
  \begin{boxedProof}
    \begin{boxFrame}
      \boxLine{Velja $\phi$:}
      \boxLine[1]{\brac{dokaz $\phi$}}
    \end{boxFrame}
  \end{boxedProof}
  \oznaka{pravilo o izključenem tretjem}
  {\footnotesize Za vsako izjavo $\phi$ velja $\phi \lor \neg\phi$.}
\end{razdelek}

\end{multicols*}

\end{document}

%%% Local Variables:
%%% mode: latex
%%% TeX-engine: luatex
%%% TeX-master: t
%%% End:
